% =============================================================================
% SIMPLIFIED PREAMBLE FOR INTRO TO MATHEMATICAL PROGRAMMING
% =============================================================================
% This is a unified, simplified preamble that consolidates all the scattered
% package loading and command definitions from the original preamble files.
%
% Original files consolidated:
%   - packages-and-commands.tex
%   - preamble-optimization.tex
%   - preamble0.tex
%   - preamble0-biblatex.tex
%   - preamble-jupyter.tex
%   - LyryxMathStyle.sty
%   - LyryxLinAlgStyle.sty
%
% Package loading order matters! This file is organized to avoid conflicts.
% =============================================================================

% =============================================================================
% SECTION 1: CORE PACKAGES (must load early)
% =============================================================================

% Language and encoding
\usepackage[english]{babel}
\usepackage[T1]{fontenc}

% Conditional statements (needed early for conditional loading)
\usepackage{ifthen}
\usepackage{iftex}
\usepackage{etoolbox}

% Font packages
\usepackage{mathptmx}  % Times-like math font
\usepackage[EULERGREEK]{sansmath}

% Simple arithmetic
\usepackage{calc}

% =============================================================================
% SECTION 2: MATH PACKAGES
% =============================================================================

\usepackage{amsmath}
\usepackage{amsfonts}
\usepackage{amsthm}
\usepackage{amssymb}
\usepackage{amscd}
\usepackage{mathtools}
\usepackage{mathrsfs}
\usepackage{gensymb}

% =============================================================================
% SECTION 3: GRAPHICS AND COLOR
% =============================================================================

\usepackage{xcolor}
\usepackage{graphicx}
\usepackage{epsfig}
\usepackage{subcaption}
\usepackage{float}
\usepackage{rotating}

% Color definitions - consolidated from all preambles
\definecolor{red}{cmyk}{0, 1.00, 0.62, 0}
\definecolor{blue}{cmyk}{1.00, 0.34, 0.0, 0.02}
\definecolor{green}{cmyk}{0.7, 0, 1.0, 0.09}
\definecolor{yellow}{cmyk}{0, 0.16, 1.0, 0}
\definecolor{gray}{cmyk}{0, 0, 0, 0.65}
\definecolor{purple}{cmyk}{0.333, 0.867, 0, 0.059}
\definecolor{ocre}{RGB}{243,102,25}

% Link and URL colors
\definecolor{urlcolor}{rgb}{0, 0.145, 0.698}
\definecolor{linkcolor}{rgb}{0.71, 0.21, 0.01}
\definecolor{citecolor}{rgb}{0.12, 0.54, 0.11}
\definecolor{linkcolour}{HTML}{013030}

% Title page and header colors
\definecolor{headersquarecolour}{HTML}{BDBDBD}
\definecolor{lyryxcolour}{HTML}{000033}
\definecolor{lyryxlogoy}{HTML}{E2970A}
\definecolor{titletextcolour}{HTML}{044155}
\definecolor{subtitletextcolour}{HTML}{F8EAC2}
\definecolor{titlemainbgcolour}{HTML}{E2AB0A}
\definecolor{titlemainboxcolour}{HTML}{F9D876}
\definecolor{sampletextcolour}{HTML}{B19545}
\definecolor{lscshdrboxcolour}{HTML}{E2AB0A}
\definecolor{lscstextcolour}{HTML}{044155}

% Section colors
\definecolor{chaptertitlecolour}{HTML}{840125}
\definecolor{sectiontitlecolour}{HTML}{0C164B}
\definecolor{subsectiontitlecolour}{HTML}{003835}
\definecolor{subsubsectiontitlecolour}{HTML}{F5AE2A}

% Standardized optimization-model setup headers (see packages-and-commands.tex).
% Color tracks the theme accent chaptertitlecolour (maroon #840125).
\makeatletter
\@ifundefined{color@chaptertitlecolour}{\definecolor{chaptertitlecolour}{HTML}{840125}}{}
\makeatother
\providecommand{\modelhead}[1]{\noindent{\bfseries\color{chaptertitlecolour}#1:}}
\providecommand{\modelsets}{\modelhead{Sets}}
\providecommand{\modelparams}{\modelhead{Parameters}}
\providecommand{\modelvars}{\modelhead{Variables}}
\providecommand{\modelmodel}{\modelhead{Model}}

% Box colors
\definecolor{theoremcolour}{HTML}{840125}
\definecolor{definitioncolour}{HTML}{0C164B}
\definecolor{examplecolour}{HTML}{ADD8E6}
\definecolor{solutioncolour}{HTML}{005F3F}
\definecolor{outcomecolour}{HTML}{E0B659}
\definecolor{resourcecolour}{HTML}{4f8f03}

% Code highlighting colors
\colorlet{shadecolor}{blue!10}
\colorlet{warning}{red!20}
\colorlet{warnline}{red}
\colorlet{information}{green!20}
\colorlet{infoline}{green}
\colorlet{codebase}{yellow!30}
\colorlet{codekeyword}{blue}
\colorlet{codecomment}{green}
\colorlet{codestring}{red}
\colorlet{notgreen}{blue!50!yellow}

% Additional colors
\definecolor{olive}{rgb}{0.3, 0.4, 0.1}
\definecolor{fore}{RGB}{249,242,215}
\definecolor{back}{RGB}{51,51,51}
\definecolor{title}{RGB}{255,0,90}
\definecolor{dgreen}{rgb}{0, 0.6, 0}
\definecolor{gold}{rgb}{1, 0.84, 0}
\definecolor{JungleGreen}{cmyk}{0.99, 0, 0.52, 0}
\definecolor{bluegreen}{cmyk}{0.85, 0, 0.33, 0}
\definecolor{RawSienna}{cmyk}{0, 0.72, 1, 0.45}
\definecolor{Magenta}{cmyk}{0, 1, 0, 0}
\definecolor{capri}{rgb}{0.0, 0.75, 1.0}
\definecolor{cherryblossompink}{rgb}{1.0, 0.72, 0.77}

% ANSI colors for code output
\definecolor{ansi-black}{HTML}{3E424D}
\definecolor{ansi-black-intense}{HTML}{282C36}
\definecolor{ansi-red}{HTML}{E75C58}
\definecolor{ansi-red-intense}{HTML}{B22B31}
\definecolor{ansi-green}{HTML}{00A250}
\definecolor{ansi-green-intense}{HTML}{007427}
\definecolor{ansi-yellow}{HTML}{DDB62B}
\definecolor{ansi-yellow-intense}{HTML}{B27D12}
\definecolor{ansi-blue}{HTML}{208FFB}
\definecolor{ansi-blue-intense}{HTML}{0065CA}
\definecolor{ansi-magenta}{HTML}{D160C4}
\definecolor{ansi-magenta-intense}{HTML}{A03196}
\definecolor{ansi-cyan}{HTML}{60C6C8}
\definecolor{ansi-cyan-intense}{HTML}{258F8F}
\definecolor{ansi-white}{HTML}{C5C1B4}
\definecolor{ansi-white-intense}{HTML}{A1A6B2}
\definecolor{ansi-default-inverse-fg}{HTML}{FFFFFF}
\definecolor{ansi-default-inverse-bg}{HTML}{000000}
\definecolor{outerrorbackground}{HTML}{FFDFDF}

% Jupyter colors
\definecolor{incolor}{HTML}{303F9F}
\definecolor{outcolor}{HTML}{D84315}
\definecolor{cellborder}{HTML}{CFCFCF}
\definecolor{cellbackground}{HTML}{F7F7F7}

% =============================================================================
% SECTION 4: TABLE PACKAGES
% =============================================================================

\usepackage{array}
\usepackage{colortbl}
\usepackage{longtable}
\usepackage{booktabs}
\usepackage{tabularx}
\usepackage{tabu}
\usepackage{multicol}
\usepackage{multirow}
\usepackage{bigstrut}
\usepackage{blkarray}

% =============================================================================
% SECTION 5: PAGE LAYOUT
% =============================================================================

% Page dimensions
\setlength{\paperheight}{11in}
\setlength{\paperwidth}{8.5in}

% Margins and geometry
\usepackage[headheight=15pt, top=1in, bottom=0.75in, outer=0.75in, inner=0.75in]{geometry}

% Line spacing and paragraph settings
\linespread{1.05}
\setlength{\parskip}{2pt plus1pt minus1pt}

% Fancy headers
\usepackage{fancyhdr}
\pagestyle{fancy}
\renewcommand{\headrulewidth}{0pt}

% Empty pages
\usepackage{emptypage}

% Margin notes
\usepackage{marginnote}

% Page colors
\usepackage{afterpage}
\usepackage{pagecolor}

% =============================================================================
% SECTION 6: TIKZ AND DRAWING
% =============================================================================

\usepackage{tikz}
\usepackage[customcolors]{hf-tikz}
\usepackage{circuitikz}
\usepackage{pgfplots}
\usepackage{pgfplotstable}
\pgfplotsset{compat=1.17}

% TikZ libraries - consolidated
\usetikzlibrary{
    arrows, arrows.meta,
    automata,
    backgrounds,
    calc,
    calendar,
    chains,
    circuits.ee.IEC,
    decorations, decorations.pathreplacing, decorations.text,
    intersections,
    matrix,
    mindmap,
    patterns,
    petri,
    positioning,
    shadows,
    shapes, shapes.geometric, shapes.callouts,
    trees
}

% TikZ styles
\tikzstyle{arrowstyle}=[scale=2.5]
\tikzstyle{directed}=[postaction={decorate,decoration={markings,
    mark=at position .65 with {\arrow[arrowstyle]{latex}}}}]
\tikzstyle{directed1}=[postaction={decorate,decoration={markings,
    mark=at position .99 with {\arrow[arrowstyle]{latex}}}}]
\tikzstyle{vertex}=[circle,fill=black!25,minimum size=20pt,inner sep=0pt]

% hf-tikz settings
\hfsetfillcolor{red!02}
\hfsetbordercolor{red}

% =============================================================================
% SECTION 7: BOXES AND FRAMES
% =============================================================================

\usepackage{framed}
\usepackage[framemethod=tikz]{mdframed}
\mdfsetup{skipabove=\topskip, skipbelow=\topskip}

\usepackage{tcolorbox}
\tcbuselibrary{theorems, breakable, skins, listingsutf8, xparse}

% Define tcbfloat type for floating tcolorboxes
\newfloat{tcbfloat}{htbp}{lot}

% mdframed styles
\mdfdefinestyle{problem}{
    backgroundcolor=shadecolor,
    skipabove=10pt,
    skipbelow=10pt,
    leftmargin=20pt,
    rightmargin=20pt,
    innertopmargin=10pt,
    innerbottommargin=10pt,
    innerleftmargin=10pt,
    middlelinewidth=0pt,
    everyline=true,
    linecolor=blue,
    linewidth=2pt
}

\newmdenv[
    skipabove=10pt,
    skipbelow=10pt,
    leftmargin=20pt,
    rightmargin=20pt,
    backgroundcolor=warning,
    innertopmargin=10pt,
    innerbottommargin=10pt,
    innerleftmargin=10pt,
    middlelinewidth=0pt,
    everyline=true,
    linecolor=warnline,
    linewidth=2pt,
    font=\normalfont\normalsize,
    frametitlefont=\large\bfseries,
    frametitleaboveskip=1em,
    frametitlerule=true,
    frametitle={\sc Caution!}
]{warn}

\newmdenv[
    skipabove=10pt,
    skipbelow=10pt,
    leftmargin=20pt,
    rightmargin=20pt,
    backgroundcolor=information,
    outerlinewidth=0pt,
    outerlinecolor=infoline,
    innertopmargin=10pt,
    innerbottommargin=10pt,
    innerleftmargin=10pt,
    middlelinewidth=0pt,
    everyline=true,
    linecolor=infoline,
    linewidth=2pt,
    font=\normalfont\normalsize,
    frametitlefont=\large\bfseries,
    frametitleaboveskip=1em,
    frametitlerule=true,
    frametitle={\sc Note}
]{info}

% tcolorbox styles
\tcbset{
    defstyle/.style={
        fonttitle=\bfseries\upshape, fontupper=\slshape,
        colback=definitioncolour!5, colframe=definitioncolour,
        before={\par\bigskip\pagebreak[0]\parindent=0pt},
        after={\par\bigskip}, top=1mm, bottom=1mm
    },
    theostyle/.style={
        fonttitle=\bfseries\upshape, fontupper=\slshape,
        colback=chaptertitlecolour!5, colframe=chaptertitlecolour,
        before={\par\bigskip\pagebreak[0]\parindent=0pt},
        after={\par\bigskip}, top=1mm, bottom=1mm
    },
    examstyle/.style={
        breakable, fonttitle=\bfseries\upshape, fontupper=\slshape,
        colback=examplecolour!5, colframe=examplecolour,
        before={\par\bigskip\pagebreak[0]\parindent=0pt},
        after={\par\bigskip}, top=1mm, bottom=1mm
    },
    solutionstyle/.style={
        fonttitle=\bfseries\upshape, fontupper=\slshape,
        colback=solutioncolour!5, colframe=solutioncolour,
        before={\par\bigskip\pagebreak[0]\parindent=0pt},
        after={\par\bigskip}, top=1mm, bottom=1mm
    }
}

% Special tcolorboxes
\newtcolorbox{titlemainbox}{
    colback=titlemainboxcolour, colframe=titlemainboxcolour,
    before={\par\bigskip\pagebreak[0]\parindent=0pt},
    after={\par\bigskip}, top=1em, bottom=1em
}

\newtcolorbox{lscshdrbox}{
    colback=lscshdrboxcolour, colframe=lscshdrboxcolour,
    top=0mm, bottom=0mm, width=19em
}

\newtcolorbox{outcome}{
    fonttitle=\bfseries\upshape, fontupper=\slshape,
    colback=white, colframe=outcomecolour,
    before={\par\bigskip\pagebreak[0]\parindent=0pt},
    after={\par\bigskip}, top=1mm, bottom=1mm, title=Outcomes
}

\newtcolorbox{resource}{
    fonttitle=\bfseries\upshape, fontupper=\slshape,
    colback=white, colframe=resourcecolour,
    before={\par\bigskip\pagebreak[0]\parindent=0pt},
    after={\par\bigskip}, top=1mm, bottom=1mm, title=Resources
}

\newtcolorbox{casestudybox}[1]{
    colback=blue!5!white, colframe=blue!75!black,
    rounded corners, boxrule=1pt, width=\textwidth,
    title=\large\textbf{Case Study:}\\ \underline{{#1}},
    breakable
}

\newtcolorbox{pioneerspotlight}[1]{
    colback=green!5!white, colframe=green!75!black,
    rounded corners, boxrule=1pt, width=\textwidth,
    title={\centering\large\textbf{Pioneer Spotlight:} \underline{#1}\par},
    breakable
}

% =============================================================================
% SECTION 8: ALGORITHMS AND CODE
% =============================================================================

% Algorithm packages - try to load, provide fallbacks if not available
\IfFileExists{algorithmicx.sty}{%
    \usepackage{algorithmicx}
    \usepackage[plain]{algorithm}
    \usepackage[noend]{algpseudocode}
    \floatstyle{ruled}
    \restylefloat{algorithm}
    % Algorithm substate environment
    \newcounter{algsubstate}
    \renewcommand{\thealgsubstate}{\alph{algsubstate}}
    \newenvironment{algsubstates}{%
        \setcounter{algsubstate}{0}%
        \renewcommand{\State}{%
            \stepcounter{algsubstate}%
            \Statex {\footnotesize\thealgsubstate:}\space}%
    }{}
    % Full line comments in Algorithmic
    \algnewcommand{\LineComment}[1]{\State \(\triangleright\) #1}
}{%
    % Fallback: provide dummy environments if algorithmicx not available
    % =============================================================================
% ALGORITHM FALLBACK DEFINITIONS
% =============================================================================
% This file provides fallback definitions for algorithm environments when
% the algorithmicx package is not available.
% =============================================================================

% Algorithm environment (float-like)
\newenvironment{algorithm}[1][]{%
    \begin{center}\begin{minipage}{0.9\textwidth}\hrule\medskip\textbf{Algorithm}\medskip
}{%
    \medskip\hrule\end{minipage}\end{center}
}

% Algorithmic environment (for pseudocode)
\newenvironment{algorithmic}[1][]{%
    \begin{quote}\ttfamily
}{%
    \end{quote}
}

% Basic commands
\newcommand{\State}{\par\noindent}
\newcommand{\Statex}{\par\noindent}

% Control flow
\newcommand{\If}[1]{\textbf{if} #1 \textbf{then}}
\newcommand{\ElsIf}[1]{\textbf{else if} #1 \textbf{then}}
\newcommand{\Else}{\textbf{else}}
\newcommand{\EndIf}{\textbf{end if}}
\newcommand{\For}[1]{\textbf{for} #1 \textbf{do}}
\newcommand{\EndFor}{\textbf{end for}}
\newcommand{\ForAll}[1]{\textbf{for all} #1 \textbf{do}}
\newcommand{\While}[1]{\textbf{while} #1 \textbf{do}}
\newcommand{\EndWhile}{\textbf{end while}}
\newcommand{\Repeat}{\textbf{repeat}}
\newcommand{\Until}[1]{\textbf{until} #1}
\newcommand{\Loop}{\textbf{loop}}
\newcommand{\EndLoop}{\textbf{end loop}}

% Return and I/O
\newcommand{\Return}[1]{\textbf{return} #1}
\newcommand{\Require}[1]{\textbf{Require:} #1}
\newcommand{\Ensure}[1]{\textbf{Ensure:} #1}
\newcommand{\Print}[1]{\textbf{print} #1}

% Functions and procedures
\newcommand{\Function}[2]{\textbf{function} \textsc{#1}(#2)}
\newcommand{\EndFunction}{\textbf{end function}}
\newcommand{\Procedure}[2]{\textbf{procedure} \textsc{#1}(#2)}
\newcommand{\EndProcedure}{\textbf{end procedure}}
\newcommand{\Call}[2]{\textsc{#1}(#2)}

% Comments
\newcommand{\Comment}[1]{\hfill\(\triangleright\) \textit{#1}}
\newcommand{\LineComment}[1]{\par\noindent\(\triangleright\) #1}

% Algorithm substate environment
\newcounter{algsubstate}
\newenvironment{algsubstates}{}{}

% =============================================================================
% END OF ALGORITHM FALLBACK
% =============================================================================

}

% Listings package for code
\usepackage{listings}

\lstset{
    language=Python,
    backgroundcolor=\color{codebase},
    tabsize=4,
    basewidth=.5em,
    rulecolor=\color{yellow},
    basicstyle=\normalsize\ttfamily,
    upquote=true,
    columns=fixed,
    extendedchars=true,
    breaklines=true,
    prebreak=\raisebox{0ex}[0ex][0ex]{\ensuremath{\hookleftarrow}},
    frame=lrtb,
    xleftmargin=5pt,
    xrightmargin=5pt,
    framesep=4pt,
    framerule=2pt,
    showtabs=false,
    showspaces=false,
    showstringspaces=false,
    morestring=[s]{"""}{"""},
    morestring=[s]{'''}{'''},
    keywordstyle=\color{codekeyword},
    commentstyle=\color{codecomment},
    stringstyle=\color{codestring},
    title=\lstname,
    captionpos=b,
    abovecaptionskip=-5pt,
    belowcaptionskip=-5pt,
    moredelim=[is][\color{black}]{<<}{>>},
    moredelim=[is][\color{red}]{<r<}{>r>},
    moredelim=[is][\color{blue}]{<b<}{>b>},
    moredelim=[is][\color{green}]{<g<}{>g>},
    moredelim=[is][\color{purple}]{<p<}{>p>},
    morekeywords={assert, bytes, self, super, with, as, yield, True, False, None}
}

\lstdefinestyle{FromFile}{
    language=Python, frame=single, numbers=left, numberstyle=\tiny,
    stepnumber=2, numbersep=7pt, numberfirstline=true,
    abovecaptionskip=2pt, belowcaptionskip=2pt
}

\lstdefinestyle{ShellOutput}{language=}
\lstdefinestyle{ShellInput}{language=}
\lstdefinestyle{pseudo}{
    basicstyle=\rmfamily, upquote=true,
    keywordstyle=\color{black}\bfseries,
    commentstyle=\color[rgb]{0.133,0.545,0.133},
    stringstyle=\color[rgb]{0.627,0.126,0.941}
}

% Inline code commands
\newcommand{\pseudoli}[1]{\lstinline[style=pseudo]!#1!}
\newcommand{\li}[1]{\lstinline[prebreak=]!#1!}
\newcommand{\lif}[1]{\lstinline[basicstyle=\footnotesize\ttfamily,language=Python,prebreak=]!#1!}
\newcommand{\lsql}[1]{\lstinline[language=SQL,prebreak=,morekeywords={TEXT, REAL, IF}]!#1!}
\newcommand{\code}[1]{\texttt{#1}}

% Verbatim and fancy verbatim
\usepackage{verbatim}
\usepackage{fancyvrb}
\usepackage{upquote}

% =============================================================================
% SECTION 9: HYPERLINKS AND REFERENCES
% =============================================================================

\usepackage{hyperref}
\hypersetup{
    pdfborder=0 0 0,
    colorlinks=true,
    allcolors=linkcolour,
    linkcolor=linkcolour,
    urlcolor=linkcolour,
    citecolor=citecolor,
    breaklinks=true
}

\usepackage{url}

% PDF border style
\makeatletter
\Hy@AtBeginDocument{
    \def\@pdfborder{0 0 1}
    \def\@pdfborderstyle{/S/U/W 1}
}
\makeatother

% =============================================================================
% SECTION 10: CAPTIONS AND FLOATS
% =============================================================================

\usepackage[format=hang, font=bf]{caption}

% Float settings
\renewcommand\floatpagefraction{.9}
\renewcommand\topfraction{.9}
\renewcommand\bottomfraction{.9}
\renewcommand\textfraction{.1}
\setcounter{totalnumber}{50}
\setcounter{topnumber}{50}
\setcounter{bottomnumber}{50}

% Make all floats centered
\makeatletter
\g@addto@macro\@floatboxreset\centering
\makeatother

% =============================================================================
% SECTION 11: SECTION HEADERS
% =============================================================================

\usepackage[explicit, sc, nobottomtitles]{titlesec}

% Chapter Header
\titleformat{\chapter}[display]
    {\color{sectiontitlecolour}\bf\Huge}
    {\hspace*{0pt}\thechapter.~#1}
    {-20pt}
    {\hspace*{0pt}{\color{sectiontitlecolour}\rule{\dimexpr\textwidth\relax}{2pt}}\Huge}
\titleformat{name=\chapter,numberless}[display]
    {\color{sectiontitlecolour}\bf\Huge}
    {\hspace*{0pt}#1}
    {-20pt}
    {\hspace*{0pt}{\color{sectiontitlecolour}\rule{\dimexpr\textwidth\relax}{2pt}}\Huge}
\titlespacing*{\chapter}{0pt}{10pt}{10pt}

% Section Header
\titleformat{\section}[display]
    {\bigskip\bigskip\hypersetup{linkcolor=sectiontitlecolour}\color{sectiontitlecolour}\bf\LARGE}
    {\hspace*{0pt}\thesection~#1}
    {-15pt}
    {\hspace*{0pt}{\color{sectiontitlecolour}\rule{\dimexpr\textwidth\relax}{2pt}}\LARGE}
\titleformat{name=\section,numberless}[display]
    {\bigskip\bigskip\hypersetup{linkcolor=sectiontitlecolour}\color{sectiontitlecolour}\bf\LARGE}
    {\hspace*{0pt}#1}
    {-15pt}
    {\hspace*{0pt}{\color{sectiontitlecolour}\rule{\dimexpr\textwidth\relax}{2pt}}\LARGE}
\titlespacing*{\section}{0pt}{0pt}{10pt}

% Subsection header
\titleformat{\subsection}[display]
    {\color{sectiontitlecolour}\bf\large}
    {\hspace*{0pt}\thesubsection.~#1}
    {-10pt}
    {\hspace*{0pt}{\color{sectiontitlecolour}\rule{\dimexpr\textwidth\relax}{2pt}}\large}
\titleformat{name=\subsection,numberless}[display]
    {\color{subsectiontitlecolour}\bf\large}
    {\hspace*{0pt}#1}
    {-10pt}
    {\hspace*{0pt}{\color{sectiontitlecolour}\rule{\dimexpr\textwidth\relax}{2pt}}\large}
\titlespacing*{\subsection}{0pt}{10pt}{10pt}

% Subsubsection header
\titleformat{\subsubsection}[display]
    {\color{sectiontitlecolour}\bf\normalsize}
    {\hspace*{0pt}\thesubsubsection.~#1}
    {-10pt}
    {\hspace*{0pt}{\color{sectiontitlecolour}\rule{\dimexpr\textwidth\relax}{2pt}}\large}
\titleformat{name=\subsubsection,numberless}[display]
    {\color{sectiontitlecolour}\bf\normalsize}
    {\hspace*{0pt}#1}
    {-10pt}
    {\hspace*{0pt}{\color{sectiontitlecolour}\rule{\dimexpr\textwidth\relax}{2pt}}\large}
\titlespacing*{\subsubsection}{0pt}{10pt}{10pt}

% Numbering depth
\setcounter{secnumdepth}{3}
\setcounter{tocdepth}{3}

% =============================================================================
% SECTION 12: INDEX AND BIBLIOGRAPHY
% =============================================================================

\usepackage{makeidx}
\makeindex

\usepackage{appendix}

% Bibliography - biblatex with fallback
\usepackage{csquotes}

\IfFileExists{biblatex.sty}{%
    \usepackage[backend=biber, giveninits=true, maxnames=100, maxcitenames=5, sorting=nyt, backref=true]{biblatex}
    \newif\ifbiblatexloaded
    \biblatexloadedtrue
    % Print url only if no doi
    \renewbibmacro*{doi+eprint+url}{%
        \printfield{doi}%
        \newunit\newblock%
        \iffieldundef{doi}{%
            \iftoggle{bbx:eprint}{\usebibmacro{eprint}}{}%
        }{}%
        \newunit\newblock%
        \iffieldundef{doi}{\usebibmacro{url+urldate}}{}%
    }
}{%
    % Biblatex not available - load fallback definitions
    % =============================================================================
% BIBLATEX FALLBACK DEFINITIONS
% =============================================================================
% This file provides fallback definitions when biblatex is not available.
% It provides dummy commands so documents can compile without biblatex.
% =============================================================================

\newif\ifbiblatexloaded
\biblatexloadedfalse

% Dummy cite commands
\makeatletter
\providecommand{\citetitle}{\@ifstar\citetitlestar\citetitlenostar}
\newcommand{\citetitlestar}[1]{[Title: #1]}
\newcommand{\citetitlenostar}[1]{[Title: #1]}
\makeatother

\providecommand{\citeauthor}[1]{[Author]}
\providecommand{\citedate}[1]{[Date]}
\providecommand{\citeurl}[1]{[URL]}

% Additional biblatex citation commands used in content
\providecommand{\parencite}[2][]{[#2]}
\providecommand{\textcite}[2][]{[#2]}
\providecommand{\autocite}[2][]{[#2]}
\providecommand{\fullcite}[1]{[Full citation: #1]}
\providecommand{\footfullcite}[1]{\footnote{[#1]}}
\providecommand{\citeyear}[1]{[Year]}
\providecommand{\citenum}[1]{[#1]}
\providecommand{\volcite}[3][]{[#2, #3]}

% Dummy print bibliography
\providecommand{\printbibliography}[1][]{}

% Dummy refsection environment
\newenvironment{refsection}[1][]{}{}

% Dummy addbibresource
\providecommand{\addbibresource}[1]{}

% =============================================================================
% END OF BIBLATEX FALLBACK
% =============================================================================

}

% =============================================================================
% SECTION 13: THEOREM ENVIRONMENTS
% =============================================================================

\theoremstyle{definition}
\newtheorem{definition}{Definition}[chapter]
\newtheorem{theorem}{Theorem}[chapter]
\newtheorem{corollary}{Corollary}[chapter]
\newtheorem{example}{Example}[chapter]
\newtheorem{proposition}{Proposition}[chapter]
\newtheorem{procedure}{Procedure}[chapter]
\newtheorem{exercise}{Exercise}[chapter]
\newtheorem{problem}{Problem}[chapter]
\newtheorem{lemma}{Lemma}[chapter]
\newtheorem{observation}{Observation}[chapter]
\newtheorem{algorithmThm}{Algorithm}[chapter]
\newtheorem{notation}{Notation}[chapter]
\newtheorem{solutionbox}{Solution}[chapter]
\newtheorem{remark}{Remark}[chapter]
\newtheorem{ICE}{ICE}

% Apply tcolorbox styling
\tcolorboxenvironment{definition}{defstyle}
\tcolorboxenvironment{theorem}{theostyle}
\tcolorboxenvironment{corollary}{theostyle}
\tcolorboxenvironment{example}{examstyle}
\tcolorboxenvironment{proposition}{theostyle}
\tcolorboxenvironment{procedure}{solutionstyle}
\tcolorboxenvironment{exercise}{solutionstyle}
\tcolorboxenvironment{problem}{solutionstyle}
\tcolorboxenvironment{lemma}{theostyle}
\tcolorboxenvironment{observation}{examstyle}
\tcolorboxenvironment{algorithmThm}{examstyle}
\tcolorboxenvironment{notation}{defstyle}
\tcolorboxenvironment{solutionbox}{solutionstyle}
\tcolorboxenvironment{remark}{examstyle}

% Solution and proof environments
\newcommand\xqed[1]{%
    \leavevmode\unskip\penalty9999 \hbox{}\nobreak\hfill
    \quad\hbox{#1}}
\def\prooffont#1{{{\bf{\textcolor{theoremcolour}{#1}}}}}
\def\solutionfont#1{{{\bf{\textcolor{solutioncolour}{#1}}}}}

\def\solution{\par\noindent{\solutionfont{\bf{\underline{Solution.}}}} \ignorespaces}
\def\endsolution{\solutionfont{\bf\xqed{$\spadesuit$}}\par\medskip}
\def\proof{\par\noindent{\prooffont{\bf{\underline{Proof.}}}} \ignorespaces}
\def\endproof{\prooffont{\bf\xqed{$\spadesuit$}}\par\medskip}

% =============================================================================
% SECTION 14: EXERCISES AND SOLUTIONS
% =============================================================================

\usepackage{answers}

% Exercise counter
\newcounter{exo}[chapter]
\newcounter{bookex}[chapter]
\renewcommand{\thebookex}{\thechapter.\arabic{bookex}}
\renewcommand{\theexo}{\thechapter.\arabic{exo}}

% Exercise environment
\newenvironment{ex}{%
    \par\medskip\noindent\refstepcounter{exo}\textbf{Exercise \theexo}\hspace{0.25em}\itshape
}{\par\medskip}

%% Unified redefinition (kept in sync with packages-and-commands.tex): accepts
%%   \begin{ex} | \begin{ex}{Title} | \begin{ex}{Title}{} | \begin{ex}[Title]{}
%% Blank titles {} are skipped (\notblank, etoolbox); second brace group ignored.
%% CAUTION: trailing brace groups are grabbed greedily, so an exercise body
%% must not BEGIN with a bare {...} group (none do as of 2026-07).
\makeatletter
\providecommand{\ex@print@title}[1]{\ \textbf{(#1)}\def\@currentlabelname{#1}}
\providecommand{\ex@peek@title}{\@ifnextchar\bgroup{\ex@grab@title}{\ex@finish}}
\providecommand{\ex@grab@title}[1]{\notblank{#1}{\ex@print@title{#1}}{}\ex@peek@second}
\providecommand{\ex@peek@second}{\@ifnextchar\bgroup{\ex@grab@second}{\ex@finish}}
\providecommand{\ex@grab@second}[1]{\ex@finish}
\providecommand{\ex@finish}{\hspace{0.25em}\itshape}
\RenewDocumentEnvironment{ex}{o}{%
  \par\medskip\noindent\refstepcounter{bookex}%
  \textbf{Exercise \thebookex}%
  \IfValueTF{#1}%
    {\notblank{#1}{\ex@print@title{#1}}{}\ex@peek@second}%
    {\ex@peek@title}%
}{%
  \par\medskip
}
\makeatother

% Solution output environment
\newenvironment{solout}[1]{%
    \par\medskip\noindent\textbf{#1}\par
}{\par\medskip}
\newcommand{\soloutparams}{{\theexo}}

% Solution association
\Newassociation{sol}{Answer}{solutions}

% =============================================================================
% SECTION 15: MISC PACKAGES
% =============================================================================

\usepackage{import}
\usepackage{pdfpages}
\usepackage{textcomp}
\usepackage{relsize}
\usepackage{scalerel}
\usepackage{paralist}
\usepackage{mdwlist}
\usepackage{eepic}
\usepackage{wasysym}
\usepackage[cc]{titlepic}
\usepackage{xparse}
\usepackage{xargs}
\usepackage{eurosym}
\usepackage{grffile}
\usepackage{adjustbox}
\usepackage{titling}
\usepackage{parskip}

% TODO notes
\usepackage[disable]{todonotes}

% =============================================================================
% SECTION 15b: BOOK INFO (needed for license page from Lyryx)
% =============================================================================

% These commands are used by the license.tex file from the Lyryx linear algebra content
\providecommand{\booktitle}{Mathematical Programming}
\providecommand{\booksubtitle}{Introduction to}
\providecommand{\bookfulltitle}{Introduction to Mathematical Programming}
\providecommand{\booksubject}{OPTIMIZATION}
\providecommand{\bookauthor}{Robert Hildebrand}
\providecommand{\version}{2025}
\providecommand{\revision}{A}

% =============================================================================
% SECTION 15c: TITLE DEFINITION
% =============================================================================

% Load title definition - works from both base and book/ directories
\IfFileExists{preamble/title.tex}{%
     \title{ \textbf{Mathematical Programming and\\ Operations Research}\\ Modeling, Algorithms, and Complexity\\
 Examples in Excel and Python}
  \author{Robert Hildebrand}

%%% Local Variables:
%%% mode: latex
%%% TeX-master: "open-optimization"
%%% End:
%
}{%
    \IfFileExists{book/preamble/title.tex}{%
         \title{ \textbf{Mathematical Programming and\\ Operations Research}\\ Modeling, Algorithms, and Complexity\\
 Examples in Excel and Python}
  \author{Robert Hildebrand}

%%% Local Variables:
%%% mode: latex
%%% TeX-master: "open-optimization"
%%% End:
%
    }{%
        % Fallback title if no title.tex found
        \title{Mathematical Programming and Operations Research}
        \author{Robert Hildebrand}
    }%
}

% =============================================================================
% SECTION 16: MATH COMMANDS
% =============================================================================

% Bold vectors
% Note: Some commands like \b, \c, \u are built-in LaTeX accent commands
% so we use \renewcommand for those. Others we define fresh.
\def\0{\mathbf{0}}
\def\a{\mathbf{a}}
\renewcommand{\b}{\mathbf{b}}  % Built-in: bar-under accent
\renewcommand{\c}{\mathbf{c}}  % Built-in: cedilla accent
\def\e{\mathbf{e}}
\def\f{\mathbf{f}}
\def\g{\mathbf{g}}
\def\p{\mathbf{p}}
\def\q{\mathbf{q}}
\renewcommand{\u}{\mathbf{u}}  % Built-in: breve accent
\def\v{\mathbf{v}}
\def\w{\mathbf{w}}
\def\x{\mathbf{x}}             % Not a built-in command
\def\y{\mathbf{y}}
\def\z{\mathbf{z}}

% Calligraphic letters
\def\CalL{\mathcal{L}}
\def\CalO{\mathcal{O}}
\def\CalV{\mathcal{V}}
\def\CalU{\mathcal{U}}
\def\PP{\mathcal{P}}
\def\NN{\mathcal{N}}
\def\II{\mathcal{I}}
\def\ZZ{\mathcal{Z}}
\def\SS{\mathcal{S}}
\def\FF{\mathcal{F}}
\def\CC{\mathcal{C}}
\newcommand{\F}{\mathcal{F}}

% Blackboard bold
\def\nn{\mathbb{N}}
\def\R{\mathbb{R}}
\def\Z{\mathbb{Z}}
\def\Q{\mathbb{Q}}
\def\cc{\mathbb{C}}
\def\one{\mathbb{1}}

% Operators
\DeclareMathOperator{\sign}{sign}
\DeclareMathOperator{\sech}{sech}
\DeclareMathOperator{\Out}{Out}
\DeclareMathOperator{\In}{In}
\DeclareMathOperator{\erf}{erf}
\DeclareMathOperator{\dom}{dom}
\DeclareMathOperator{\cl}{cl}
\DeclareMathOperator{\intr}{intr}
\DeclareMathOperator{\argmin}{arg\,min}
\newcommand{\argmax}{\mbox{argmax}}

% Text operators
\def\rank{\textup{rank}}
\def\conv{\textup{conv}}
\def\cone{\textup{cone}}
\def\supp{\textup{supp}}
\def\rc{\textup{rec.cone}}
\def\ri{\textup{rel.int}}
\def\rb{\textup{rel.bd}}
\def\bd{\textup{bd}}
\def\ls{\textup{lin.space}}
\def\aff{\textup{aff}}
\def\st{\text{ s.t. }}

% Delimiters
\providecommand{\abs}[1]{\left\lvert#1\right\rvert}
\providecommand{\norm}[1]{\left\lVert#1\right\rVert}
\providecommand{\set}[1]{\lbrace#1\rbrace}
\providecommand{\setconstruct}[2]{\lbrace#1:#2\rbrace}
\newcommand{\floor}[1]{\lfloor #1 \rfloor}
\newcommand{\ceil}[1]{\lceil #1 \rceil}

% Matrix and vector
\newcommand{\trp}{^{\mathsf T}}
\newcommand{\hrm}{^{\mathsf H}}
\newcommand{\T}{\mathsf{T}}
\newcommand{\ipt}[2]{\langle #1,#2 \rangle}
\newcommand{\ip}{\int_{-\infty}^{+\infty}}
\renewcommand{\ker}[1]{\mathcal{N}(#1)}
\newcommand{\ran}[1]{\mathcal{R}(#1)}
\renewcommand{\vec}[1]{\mathbf{#1}}
\newcommand{\mm}[1]{\mathbf{#1}}
\newcommand{\vect}{\vec}
\newcommand{\longvect}{\overrightarrow}
\newcommand{\leftB}{\left[}
\newcommand{\rightB}{\right]}
\setcounter{MaxMatrixCols}{10}

% Epsilon
\renewcommand{\epsilon}{\varepsilon}

% Indicator function
\newcommand{\indicator}{\boldsymbol{1}}

% Misc math
\def\subspace{\lhd}
\def\bU{{\bar{u}}}
\def\M{M_n}
\def\tq{{\,:\,}}
\newcommand{\ssep}{~:~}
\def\lt{<}
\def\gt{>}

% For block arrays
\newcommand\topstrut[1][0.8ex]{\setlength\bigstrutjot{#1}{\bigstrut[t]}}
\newcommand\botstrut[1][0.6ex]{\setlength\bigstrutjot{#1}{\bigstrut[b]}}
\def\lvl#1{\multicolumn{1}{|c}{#1}}
\def\rvl#1{\multicolumn{1}{c|}{#1}}

% =============================================================================
% SECTION 17: PSEUDOCODE COMMANDS
% =============================================================================

\NewDocumentCommand\allocate{m+g}{%
    \IfNoValueTF{#2}
        {\mathrm{empty}(#1)}
        {\mathrm{empty}(#1, #2)}
}

\NewDocumentCommand\zeros{m+g}{%
    \IfNoValueTF{#2}
        {\mathrm{zeros}(#1)}
        {\mathrm{zeros}(#1, #2)}%
}

\newcommand{\Id}[1]{\mathrm{Id}(#1)}
\newcommand{\makecopy}[1]{\mathrm{copy}(#1)}
\newcommand{\shape}[1]{\mathrm{shape}(#1)}
\newcommand{\size}[1]{\mathrm{size}(#1)}

% =============================================================================
% SECTION 18: COMPLEXITY AND SPECIAL NOTATION
% =============================================================================

\newcommand{\complexity}[1]{{\color{red} \textit{#1}}}
\newcommand{\nphard}{\complexity{NP-Hard}}
\newcommand{\npcomplete}{\complexity{NP-Complete}}
\newcommand{\polynomial}{{\color{ansi-green} \textit{Polynomial time (P)}}}
\newcommand{\true}{\tgreen{\texttt{ true }}}
\newcommand{\false}{\tred{\texttt{ false }}}
\newcommand{\NP}{\text{NP}}
\newcommand{\coNP}{\text{coNP}}
\def\BQP{\mathrm{BQP}}
\def\xLP{x^{\mathrm{LP}}}

% Color text commands
\newcommand{\tblue}[1]{{\color{blue} #1}}
\newcommand{\tred}[1]{{\color{red} #1}}
\newcommand{\tgreen}[1]{{\color{notgreen} #1}}

% Highlight cell in tables
\newcommand{\hi}{\cellcolor{red}}

% =============================================================================
% SECTION 19: SPECIAL ENVIRONMENTS
% =============================================================================

% Lab environment
\renewcommand{\chaptername}{Lab}
\newcommand{\lab}[2]{\chapter[#2]{#1}}
\newcommand{\objective}[1]{{\bf Lab Objective: } \emph{#1} \bigskip}

% General framed environment
% Note: These environments use @ in internal commands, so we need makeatletter
\makeatletter
\newcounter{general}
\newenvironment{general}[2]{%
    \par\medskip
    \protected@edef\@currentlabelname{#1}%
    \begin{framed}
    \begingroup\color{black}%
    \textbf{#1: }\ignorespaces
    \par #2
    \par
}{\endgroup\end{framed}\medskip}

% Example with code environments
\newenvironment{optexample}{%
    \par\medskip
    \refstepcounter{exo}%
    \protected@edef\@currentlabelname{Example \theexo}%
    \begin{framed}
    \begingroup\color{blue}%
    \textbf{Example \theexo: }\ignorespaces
}{\endgroup\end{framed}\medskip}

\newenvironment{examplewithcode}[2]{%
    \par\medskip
    \refstepcounter{exo}%
    \protected@edef\@currentlabelname{Example \theexo}%
    \begingroup
    \begin{center}
    \rule{\textwidth}{0.4pt}\\
    \vspace{-0.2cm}
    \rule{\textwidth}{0.4pt}
    \end{center}
    \color{black}%
    \textbf{Example \theexo: #1}%
    \hfill \protect{#2}%
    \ignorespaces
    \par
}{%
    \endgroup
    \begin{center}
    \rule{\textwidth}{1.8pt}
    \end{center}
    \medskip
}

\newenvironment{examplewithoutcode}[2]{%
    \par\medskip
    \refstepcounter{exo}%
    \protected@edef\@currentlabelname{Example \theexo}%
    \begin{framed}
    \begingroup\color{black}%
    \textbf{Example \theexo: #1}%
    \ignorespaces
    \par
}{\endgroup\end{framed}\medskip}

\newenvironment{examplewithallcode}[4]{%
    \par\medskip
    \refstepcounter{exo}%
    \protected@edef\@currentlabelname{Example \theexo}%
    \begin{framed}
    \begingroup\color{black}%
    \textbf{Example \theexo: #1}%
    \hfill \protect\href{#2}{\ Excel \ }\protect\href{#3}{\ PuLP \ }\protect\href{#4}{\ Gurobipy \ }%
    \ignorespaces
    \par
}{\endgroup\end{framed}\medskip}
\makeatother

% Contributors environment
\newcommand{\contributor}[2]{\noindent\vtop{\hsize14pc#1\\\itshape#2}\par}
\newenvironment{contributors}{{List of Contributors}%
    \thispagestyle{plain}\begin{multicols}{2}\parindent0pt%
    \parskip6pt plus2pt minus1pt%
    \widowpenalty10000\clubpenalty10000}
{\end{multicols}
This project is funded in part by the National Science Foundation, grant no. TUES Phase II DUE-1323785.
\clearpage}

% Enumerate style
\newenvironment{enumialphparenastyle}{%
    \let\oldlabelenumi=\labelenumi
    \renewcommand{\labelenumi}{\textrm{(\alph{enumi})}}
}{\let\labelenumi=\oldlabelenumi}

% Center without spacing
\newenvironment{center*}[0]{\bgroup \centering}{\egroup}

% =============================================================================
% SECTION 20: TODO COMMANDS
% =============================================================================

\newcommand\getcurrentref[1]{%
    \ifnumequal{\value{#1}}{0}
    {??}
    {\the\value{#1}}%
}

\newcommand{\todored}[1]{}

\let\Chaptermark\chaptermark
\def\chaptermark#1{\def\Chaptername{#1}\Chaptermark{#1}}
\let\Sectionmark\sectionmark
\def\sectionmark#1{\def\Sectionname{#1}\Sectionmark{#1}}
\let\Subsectionmark\subsectionmark
\def\subsectionmark#1{\def\Subsectionname{#1}\Subsectionmark{#1}}
\let\Subsubsectionmark\subsubsectionmark
\def\subsubsectionmark#1{\def\Subsubsectionname{#1}\Subsubsectionmark{#1}}

\newcommand{\todoChapter}[1]{}
\newcommand{\todoSection}[1]{}

% =============================================================================
% SECTION 21: GRAPHICS PATHS
% =============================================================================

\graphicspath{%
    {./images/},%
    {images/},%
    {Figures/},%
    {../Figures/},%
    {figures/},%
    {graph-theory-graphics/},%
    {../graph-theory-graphics/},%
    {optimization/figures/figures-static/},%
    {optimization/figures/figures-source/},%
    {optimization/multi-objective/images/},%
    {external-sources/aFirstCourseLinearAlgebra/},%
    {external-sources/lineqlpbook/},%
    {external-sources/Christopher_Griffin_Penn_State_University/},%
    {external-sources/OR-Book-Vorwerk/images/},%
    {external-sources/foundationsAppliedMathematicsLabs/},%
    {external-sources/foundationsAppliedMathematicsLabs/Volume2/Simplex/figures/},%
    {../external-sources/aFirstCourseLinearAlgebra/},%
    {../external-sources/lineqlpbook/},%
    {../external-sources/Christopher_Griffin_Penn_State_University/},%
    {../external-sources/OR-Book-Vorwerk/images/},%
    {../external-sources/foundationsAppliedMathematicsLabs/},%
    {../external-sources/foundationsAppliedMathematicsLabs/Volume2/Simplex/figures/},%
    {part1-linear-programming/figures/},%
}

% =============================================================================
% SECTION 22: LOGO AND BRAND COMMANDS
% =============================================================================

\newcommand{\julia}{\raisebox{-.2\height}{
\includegraphics[scale=0.07]{julia-logo}}}
\newcommand{\jupyter}{\raisebox{-.2\height}{
\includegraphics[scale=0.09]{jupyter-logo}}}
\newcommand{\jump}{\raisebox{-.2\height}{
\includegraphics[scale=0.08]{jump-logo}}}
\newcommand{\gurobi}{\raisebox{-.2\height}{\includegraphics[scale=0.06]{gurobi-logo}}Gurobi}
\newcommand{\coin}{\raisebox{-.2\height}{
\includegraphics[scale=0.2]{coin-or-logo}}}
\newcommand{\ipopt}{\raisebox{-.2\height}{\texttt{Ipopt}}}

% =============================================================================
% SECTION 23: FIGURE INCLUDE COMMANDS
% =============================================================================

\DeclareRobustCommand\footnotetextmetadata[1]{%
    \footnotetext{\citetitle*{#1}, from \citeurl{#1}.
    \citeauthor{#1}, \citedate{#1}.
    }
}

\newcommandx{\refincludefigurestatic}[4][1=\relax, 2={width=.4\linewidth}, 3=t, usedefault]{%
    \begin{figure}[#3]%
    \centering%
    \includegraphics[#2]{optimization/figures/figures-static/#4}\\
    \centerline{\scriptsize \copyright~\citeauthor{#4}\protect\footnotemark}
    \caption{%
        \ifx#1\relax \citetitle*{#4}\else #1\fi%
    }%
    \label{fig:#4}%
    \end{figure}%
    \autoref{fig:#4}\footnotetextmetadata{#4}%
}

\newcommandx{\includegraphicstatic}[2][1={width=.4\linewidth}, usedefault]{%
    \begin{center}
    \includegraphics[#1]{optimization/figures/figures-static/#2}\\
    \end{center}
    \centerline{\scriptsize \copyright~\citeauthor{#2}\protect\footnotemark}
    \footnotetextmetadata{#2}%
}

\newcommandx{\includegraphicsource}[2][1={width=.4\linewidth}, usedefault]{%
    \begin{center}
    \includegraphics[#1]{optimization/figures/figures-source/#2}\\
    \end{center}
    \centerline{\scriptsize \copyright~\citeauthor{#2}\protect\footnotemark}
    \footnotetextmetadata{#2}%
}

\newcommandx{\includefigurestatic}[4][1=\relax, 2={width=.4\linewidth}, 3=t, usedefault]{%
    \begin{figure}[#3]%
    \centering%
    \includegraphics[#2]{optimization/figures/figures-static/#4}\\
    \centerline{\scriptsize \copyright~\citeauthor{#4}\protect\footnotemark}
    \caption{%
        \ifx#1\relax \citetitle*{#4}\else #1\fi%
    }%
    \label{fig:#4}%
    \end{figure}%
    \footnotetextmetadata{#4}%
}

\newcommandx{\includefiguresource}[4][1=\relax, 2={width=.4\linewidth}, 3=t, usedefault]{%
    \begin{figure}[#3]%
    \centering%
    \includegraphics[#2]{optimization/figures/figures-source/#4}\\
    \centerline{\scriptsize \copyright~\citeauthor{#4}\protect\footnotemark}
    \caption{%
        \ifx#1\relax \citetitle*{#4}\else #1\fi%
    }%
    \label{fig:#4}%
    \end{figure}%
    \footnotetextmetadata{#4}%
}

% =============================================================================
% SECTION 24: LINEARE ALGEBRA SPECIFIC
% =============================================================================

% RREF notation
\newcommand{\RREF}{Reduced Row-Echelon Form}
\newcommand{\rref}{reduced row-echelon form}
\newcommand{\EF}{Row-Echelon Form}
\newcommand{\ef}{row-echelon form}

% Eigenvalue variable
\newcommand{\eigenVar}{x}

% Dot product
\newcommand{\dotprod}{\bullet}
\newcommand{\vectlength}{\|}

% Shorthand
\def\ds{\displaystyle}
\def\dfont#1{{\bf{#1}}}
\def\ifont#1{{\it{{#1}}}}
\def\xrefn#1{\ref{#1}}
\def\figrdef#1{\label{#1}}

% Math operators
\newcommand{\opft}{\mathgroup\symoperators}
\newcommand{\lip}{\mathop{\opft Lip}\nolimits}
\newcommand{\Res}{\mathop{\opft Res}\nolimits}
\newcommand{\INT}{\mathop{\opft int}\nolimits}
\newcommand{\spt}{\mathop{\opft spt}\nolimits}
\newcommand{\ann}{\mathop{\opft ann}\nolimits}
\newcommand{\diam}{\mathop{\opft diam}\nolimits}
\newcommand{\dist}{\mathop{\opft dist}\nolimits}
\providecommand{\text}[1]{\mbox{#1}}
\providecommand{\func}[1]{\mbox{#1}}
\providecommand{\limfunc}[1]{\mbox{#1}}

% Lift and upper bound
\newcommand{\lift}{\mathrm{lift}}
\newcommand{\ub}{\mathrm{ubbest}}
\newcommand{\mathup}[1]{#1}

% =============================================================================
% SECTION 25: PANDOC / JUPYTER COMPATIBILITY
% =============================================================================

\providecommand{\tightlist}{%
    \setlength{\itemsep}{0pt}\setlength{\parskip}{0pt}
}

\DefineVerbatimEnvironment{Highlighting}{Verbatim}{commandchars=\\\{\}}
\newenvironment{Shaded}{}{}

% Syntax highlighting tokens
\newcommand{\KeywordTok}[1]{\textcolor[rgb]{0.00,0.44,0.13}{\textbf{{#1}}}}
\newcommand{\DataTypeTok}[1]{\textcolor[rgb]{0.56,0.13,0.00}{{#1}}}
\newcommand{\DecValTok}[1]{\textcolor[rgb]{0.25,0.63,0.44}{{#1}}}
\newcommand{\BaseNTok}[1]{\textcolor[rgb]{0.25,0.63,0.44}{{#1}}}
\newcommand{\FloatTok}[1]{\textcolor[rgb]{0.25,0.63,0.44}{{#1}}}
\newcommand{\CharTok}[1]{\textcolor[rgb]{0.25,0.44,0.63}{{#1}}}
\newcommand{\StringTok}[1]{\textcolor[rgb]{0.25,0.44,0.63}{{#1}}}
\newcommand{\CommentTok}[1]{\textcolor[rgb]{0.38,0.63,0.69}{\textit{{#1}}}}
\newcommand{\OtherTok}[1]{\textcolor[rgb]{0.00,0.44,0.13}{{#1}}}
\newcommand{\AlertTok}[1]{\textcolor[rgb]{1.00,0.00,0.00}{\textbf{{#1}}}}
\newcommand{\FunctionTok}[1]{\textcolor[rgb]{0.02,0.16,0.49}{{#1}}}
\newcommand{\RegionMarkerTok}[1]{{#1}}
\newcommand{\ErrorTok}[1]{\textcolor[rgb]{1.00,0.00,0.00}{\textbf{{#1}}}}
\newcommand{\NormalTok}[1]{{#1}}
\newcommand{\ConstantTok}[1]{\textcolor[rgb]{0.53,0.00,0.00}{{#1}}}
\newcommand{\SpecialCharTok}[1]{\textcolor[rgb]{0.25,0.44,0.63}{{#1}}}
\newcommand{\VerbatimStringTok}[1]{\textcolor[rgb]{0.25,0.44,0.63}{{#1}}}
\newcommand{\SpecialStringTok}[1]{\textcolor[rgb]{0.73,0.40,0.53}{{#1}}}
\newcommand{\ImportTok}[1]{{#1}}
\newcommand{\DocumentationTok}[1]{\textcolor[rgb]{0.73,0.13,0.13}{\textit{{#1}}}}
\newcommand{\AnnotationTok}[1]{\textcolor[rgb]{0.38,0.63,0.69}{\textbf{\textit{{#1}}}}}
\newcommand{\CommentVarTok}[1]{\textcolor[rgb]{0.38,0.63,0.69}{\textbf{\textit{{#1}}}}}
\newcommand{\VariableTok}[1]{\textcolor[rgb]{0.10,0.09,0.49}{{#1}}}
\newcommand{\ControlFlowTok}[1]{\textcolor[rgb]{0.00,0.44,0.13}{\textbf{{#1}}}}
\newcommand{\OperatorTok}[1]{\textcolor[rgb]{0.40,0.40,0.40}{{#1}}}
\newcommand{\BuiltInTok}[1]{{#1}}
\newcommand{\ExtensionTok}[1]{{#1}}
\newcommand{\PreprocessorTok}[1]{\textcolor[rgb]{0.74,0.48,0.00}{{#1}}}
\newcommand{\AttributeTok}[1]{\textcolor[rgb]{0.49,0.56,0.16}{{#1}}}
\newcommand{\InformationTok}[1]{\textcolor[rgb]{0.38,0.63,0.69}{\textbf{\textit{{#1}}}}}
\newcommand{\WarningTok}[1]{\textcolor[rgb]{0.38,0.63,0.69}{\textbf{\textit{{#1}}}}}

% Line break command
\def\br{\hspace*{\fill} \\* }

% Jupyter prompt
\makeatletter
\newcommand{\boxspacing}{\kern\kvtcb@left@rule\kern\kvtcb@boxsep}
\makeatother
\newcommand{\prompt}[4]{%
    {\ttfamily\llap{{\color{#2}[#3]:\hspace{3pt}#4}}\vspace{-\baselineskip}}
}

% grffile fix for old versions
\makeatletter
\@ifpackagelater{grffile}{2019/11/01}{}
{
    \def\Gread@@xetex#1{%
        \IfFileExists{"\Gin@base".bb}%
        {\Gread@eps{\Gin@base.bb}}%
        {\Gread@@xetex@aux#1}%
    }
}
\makeatother

% Quotation fix
\AtBeginDocument{%
    \def\PYZsq{\textquotesingle}%
}

% Sloppy for long URLs
\sloppy

% =============================================================================
% SECTION 26: FINAL SETTINGS
% =============================================================================

% Reset footnote counter every chapter
\makeatletter
\@addtoreset{footnote}{chapter}
\makeatother

% Equation numbering
\numberwithin{equation}{section}
\renewcommand{\theequation}{\arabic{chapter}.\arabic{equation}}

% Gray square for headers
\newcommand\graysquare{\textcolor{headersquarecolour}{\rule{1ex}{1ex}}}

% Fancy headers - simplified
\fancyhf{}
\fancyhead[RO]{\rightmark\hspace{0.5em}\graysquare\hspace{0.5em}\thepage}
\fancyhead[LE]{\thepage\hspace{0.5em}\graysquare\hspace{0.5em}\leftmark}

% Cancel command
\newsavebox\CBox
\newcommand\hcancel[2][0.5pt]{%
    \ifmmode\sbox\CBox{$#2$}\else\sbox\CBox{#2}\fi%
    \makebox[0pt][l]{\usebox\CBox}%
    \rule[0.5\ht\CBox-#1/2]{\wd\CBox}{#1}
}

% Code cell counter
\newcounter{codeCell}
\makeatletter
\newenvironment{codeCell}{%
    \refstepcounter{codeCell}%
    \protected@edef\@currentlabelname{Code}%
}{}
\makeatother

% Left/right page clearing
\newcommand*\cleartoleftpage{%
    \clearpage
    \ifodd\value{page}\hbox{}\newpage\fi
}

% TikZ component node style
\tikzset{
    componentnode/.style={
        font=\scriptsize,
        align=center,
        white,
        execute at begin node=\setlength{\baselineskip}{1.25em}
    }
}

% PDF page height fix
\pdfpageheight=\paperheight

% =============================================================================
% END OF SIMPLIFIED PREAMBLE
% =============================================================================
